\section{Cartan 联络}\label{sec:cartan-connection}

设 $(x^0,x^1,x^2,x^3)=(t,x,y,z)$ 为 Galileo 时空上的惯性参考系, 因而可将该时空与 $\mathbb{R}^4$ 等同. 回顾, Newton 引力由\textbf{引力势} $\Phi:\mathbb{R}^4\to\mathbb{R}$ 描述. 这个势通过
\[
  \frac{\mathrm{d}^2x^i}{\mathrm{d}t^2}=-\frac{\partial\Phi}{\partial x^i}
\]
($i=1,2,3$) 确定自由下落粒子的运动, 而势本身又由\textbf{物质密度函数} $\rho:\mathbb{R}^4\to\mathbb{R}$ 通过 \textbf{Poisson 方程}
\[
  \frac{\partial^2\Phi}{\partial x^2}+\frac{\partial^2\Phi}{\partial y^2}
   +\frac{\partial^2\Phi}{\partial z^2}=4\pi\rho
\]
确定 (这里采用使 Newton 万有引力常数 $G$ 等于 $1$ 的单位). 真空 Poisson 方程 (对应于全部物质都集中于引力势奇点的情形) 就是著名的 \textbf{Laplace 方程}
\[
  \frac{\partial^2\Phi}{\partial x^2}+\frac{\partial^2\Phi}{\partial y^2}
   +\frac{\partial^2\Phi}{\partial z^2}=0.
\]

注意, 所有粒子的运动方程相同, 与其质量无关. 这一观察可追溯到 Galileo, Einstein 将其表述为所谓的\textbf{等效原理}. 它意味着, 引力场在 Galileo 时空上确定了特殊的曲线, 即自由下落粒子的运动曲线. 这些曲线是一个对称联络的测地线, 此联络称为 \textbf{Cartan 联络}, 由非零 Christoffel 符号
\[
  \Gamma_{00}^i=\frac{\partial\Phi}{\partial x^i}\qquad(i=1,2,3)
\]
定义 [参见 Exercise~\ref{ex:cartan-connection}(\ref{ex:cartan-geodesics-as-free-fall})], 对应的非零联络形式为
\[
  \omega_0^i=\frac{\partial\Phi}{\partial x^i}\,\mathrm{d}t.
\]
容易验证, Cartan 结构方程
\[
  \Omega_\nu^\mu=\mathrm{d}\omega_\nu^\mu+\sum_{\alpha=0}^3\omega_\alpha^\mu\wedge\omega_\nu^\alpha
\]
对任意对称联络仍成立, 因而其非零曲率形式为
\[
  \Omega_0^i=\sum_{j=1}^3\frac{\partial^2\Phi}{\partial x^j\partial x^i}\,
                    \mathrm{d}x^j\wedge\mathrm{d}t.
\]
这一联络的 Ricci 曲率张量为
\[
  \operatorname{Ric}=\left(\frac{\partial^2\Phi}{\partial x^2}
    +\frac{\partial^2\Phi}{\partial y^2}+\frac{\partial^2\Phi}{\partial z^2}\right)
      \mathrm{d}t\otimes\mathrm{d}t
\]
[参见 Exercise~\ref{ex:cartan-connection}(\ref{ex:cartan-connection-ricci-curvature})], 从而 Poisson 方程可以写为
\[
  \operatorname{Ric}=4\pi\rho\,\mathrm{d}t\otimes\mathrm{d}t.
\]
特别地, Laplace 方程可以写为
\[
  \operatorname{Ric}=0.
\]

\begin{RiemannExercise}\label{ex:cartan-connection}
\begin{enumerate}[label=(\arabic*),ref=\arabic*,leftmargin=*]
\item\label{ex:cartan-geodesics-as-free-fall} 验证自由下落粒子的运动确实是 Cartan 联络的测地线. 还有哪些其他测地线? 如何解释它们?
\item\label{ex:cartan-connection-ricci-curvature} 验证 Cartan 联络的 Ricci 曲率张量的公式.
\item\label{ex:cartan-connection-compatibility-with-galileo-structure} 证明 Cartan 联络 $\nabla$ 与 Galileo 结构相容, 即证明:
\begin{enumerate}[label=(\alph*),leftmargin=*]
\item 对所有 $X\in\mathfrak{X}(\mathbb{R}^4)$, $\nabla_X\mathrm{d}t=0$ [参见第~\ref{chap:riemannian-manifolds}~章的 Exercise~\ref{ex:affine-connections}(\ref{ex:covariant-derivative-of-tensors})];
\item 若 $E,F\in\mathfrak{X}(\mathbb{R}^4)$ 与同时超曲面相切, 且沿某条曲线 $c:\mathbb{R}\to\mathbb{R}^4$ 平行, 则 $\langle E,F\rangle$ 恒定.
\end{enumerate}
\item\label{ex:nonflat-cartan-connection-is-not-levi-civita}
证明若 Cartan 联络的曲率不为零, 则它不是 $\mathbb{R}^4$ 上任何伪 Riemann 度量的 Levi–Civita 联络 (参见第~\RGref{sec:general-relativity}~节).
\end{enumerate}
\end{RiemannExercise}
